\section{The general separator and exactification of one separator}
\label{subsec:sparse}

This appendix proves the separator with a dot-product dimension \(K_0\)
and a phase slope \(c\), and compares three ways of turning a single
separator into a restriction: sparse interpolation, direct copying, and,
for its tensor square, an average over a prime number of roots of unity.
The construction of Section~\ref{sec:extraction} uses none of this: it
needs only the case \(K_0=M\), \(c=1\) (Proposition~\ref{prop:separator}),
and it exactifies the whole catalytic degeneration of
Proposition~\ref{prop:catalytic} by Lemma~\ref{lem:exactification}.

\subsection{The separator with general parameters}

Let \(T=\sum_{g=0}^{M-1}T_g\) be a tensor with \(M\) sectors sharing the
first leg \(X\), with coordinates \(x_\alpha\), \(y_{g,\beta}\),
\(z_{g,\gamma}\) and coefficients \(T_g[\alpha,\beta,\gamma]\) as in the
proof of Proposition~\ref{prop:separator}. For positive integers \(K_0\)
and \(c\), put
\begin{equation}\label{eq:separator-parameters}
 L_0=K_0+c(M-1),\qquad
 \delta_*=\min\Bigl(M-1,\Bigl\lfloor\frac{K_0-1}{c}\Bigr\rfloor\Bigr).
\end{equation}
For \(K_0=M\) and \(c=1\) these are \(L_0=2M-1\) and \(\delta_*=M-1\), and
the next proposition is Proposition~\ref{prop:separator}.

\begin{proposition}[General separator]\label{prop:separator-general}
Let \(T\), \(K_0\), \(c\), \(L_0\) and \(\delta_*\) be as above. The
substitutions in the proof, whose entries are Laurent polynomials in
\(t\), transform \(L_0T\) into a polynomial \(\widetilde T(t)\) whose
exponents lie in \(\{c\delta^2:0\le\delta\le\delta_*\}\) and whose
constant term is \(\widetilde T_0=\bigoplus_gT_gB_X(K_0)\). In particular,
\begin{equation}\label{eq:separator-general}
 L_0T\degen\bigoplus_{g=0}^{M-1}T_gB_X(K_0),
\end{equation}
with degree at most \(c\delta_*^2\). After a permutation of the legs, the
same holds when the shared leg is the second or the third leg.
\end{proposition}

\begin{proof}
We follow the proof of Proposition~\ref{prop:separator} with
\(\zeta=\zeta_{L_0}\), source copies \(\ell=0,\ldots,L_0-1\), and target
variables \(X_{\alpha,g'}\), \(Y_{g,\beta,e}\), \(Z_{g,\gamma,e'}\) with
\(0\le g'\le M-1\) and \(0\le e,e'\le K_0-1\). In copy \(\ell\) substitute
\[
 x_\alpha\mapsto\sum_{g'}\zeta^{\ell cg'}t^{cg'^2}X_{\alpha,g'},\quad
 y_{g,\beta}\mapsto\sum_{e}\zeta^{\ell(e-cg)}t^{2ge-cg^2}Y_{g,\beta,e},\quad
 z_{g,\gamma}\mapsto\frac1{L_0}\sum_{e'}\zeta^{-\ell e'}t^{-2ge'}
 Z_{g,\gamma,e'}.
\]
Summing over \(\ell\), the term \(x_\alpha y_{g,\beta}z_{g,\gamma}\)
contributes to the coefficient of \(X_{\alpha,g'}Y_{g,\beta,e}Z_{g,\gamma,e'}\)
a power of \(t\) times the average over \(\ell\) of
\(\zeta^{\ell(e-e'+c(g'-g))}\), which is \(1\) if \(L_0\) divides
\(e-e'+c(g'-g)\) and \(0\) otherwise. Since
\(|e-e'+c(g'-g)|\le(K_0-1)+c(M-1)=L_0-1\), the surviving terms are exactly
those with \(e-e'=c(g-g')\). On them the exponent of \(t\) is
\(cg'^2+2g(e-e')-cg^2=cg'^2+2cg(g-g')-cg^2=c(g'-g)^2\). Hence the sum
of the images of the \(L_0\) copies is the polynomial
\[
 \widetilde T(t)=\sum_{\substack{g,g',e,e'\\e-e'+c(g'-g)=0}}t^{c(g'-g)^2}
 \sum_{\alpha,\beta,\gamma}T_g[\alpha,\beta,\gamma]\,
 X_{\alpha,g'}Y_{g,\beta,e}Z_{g,\gamma,e'}.
\]
The exponent vanishes exactly when \(g'=g\), and then \(e=e'\); as in
Step~3 of the proof of Proposition~\ref{prop:separator}, these terms form
\(\bigoplus_gT_gB_X(K_0)\). For every surviving term, \(\delta=|g'-g|\)
satisfies \(\delta\le M-1\) and \(c\delta=|e-e'|\le K_0-1\), so \(\delta\le\delta_*\).
Permuting the legs permutes the three substitution maps.
\end{proof}

The two parameters trade source copies against sparsity of the exponent
support: increasing \(c\) by one adds \(M-1\) source copies and never
increases the number \(\delta_*+1\) of possible exponents, which drops
below \(M\) once \(c(M-1)>K_0-1\).

\subsection{Sparse exactification and direct copying}

With \(k=1\), Lemma~\ref{lem:exactification} uses \(D+1\) evaluations for
a polynomial of degree at most \(D\). When the polynomial has few distinct
exponents, fewer evaluations suffice.

\begin{lemma}[Sparse exactification]\label{lem:sparse-exactification}
Suppose that a triple \(\Psi(t)\) of linear maps with Laurent-polynomial
entries transforms a tensor into
\(\widetilde T(t)=\sum_{\eta\in\mathcal E}t^\eta\widetilde T_\eta\), where
\(\mathcal E\) is a finite set of nonnegative integers that contains \(0\).
Then the direct sum of \(|\mathcal E|\) copies of that tensor restricts to
\(\widetilde T_0\), with rational interpolation coefficients.
\end{lemma}

\begin{proof}
The polynomial
\(\psi(u)=\prod_{\eta\in\mathcal E\setminus\{0\}}(u-2^\eta)/(1-2^\eta)
=\sum_{\iota=0}^{|\mathcal E|-1}\psi_\iota u^\iota\) has rational
coefficients, \(\psi(1)=1\), and \(\psi(2^\eta)=0\) for every
\(\eta\in\mathcal E\setminus\{0\}\). Hence
\(\sum_\iota\psi_\iota\widetilde T(2^\iota)
=\sum_{\eta\in\mathcal E}\psi(2^\eta)\widetilde T_\eta=\widetilde T_0\).
Each \(\widetilde T(2^\iota)\) is the image of one copy of the source
tensor under \(\Psi(2^\iota)\), which is defined because \(2^\iota\ne0\).
Sending the \(|\mathcal E|\) copies to the common target variables and
absorbing \(\psi_\iota\) into one leg of copy \(\iota\) gives the
restriction.
\end{proof}

For a separator the source is \(L_0T\), and by
Proposition~\ref{prop:separator-general} the exponent support has
\(\delta_*+1\) elements. Lemma~\ref{lem:sparse-exactification} therefore
gives a restriction from \(L_0(\delta_*+1)\) copies of \(T\) to
\(\bigoplus_gT_gB_X(K_0)\).

Direct copying reaches the same target from \(MK_0\) copies of \(T\) by a
restriction, with no parameter \(t\). Take one copy of \(T\) for each pair
\((g,e)\) with \(0\le g\le M-1\) and \(0\le e\le K_0-1\), restrict it to
sector \(g\) by setting the coordinates of the other sectors to zero, and
map \(x_\alpha\mapsto X_{\alpha,g}\), \(y_{g,\beta}\mapsto Y_{g,\beta,e}\)
and \(z_{g,\gamma}\mapsto Z_{g,\gamma,e}\). For fixed \(g\), the \(K_0\)
copies keep separate second and third legs but share the first leg, so
together they form \(T_gB_X(K_0)\); different values of \(g\) use
different variables on all three legs. This gives
\(MK_0\,T\restr\bigoplus_gT_gB_X(K_0)\).

Sparse exactification of one separator never uses fewer copies than direct
copying:
\begin{equation}\label{eq:one-power-limit}
 L_0(\delta_*+1)\ge MK_0.
\end{equation}
If \(\delta_*=M-1\), this follows from \(L_0\ge K_0\). Otherwise
\(\delta_*=\lfloor(K_0-1)/c\rfloor\), so \(\delta_*+1=\lceil K_0/c\rceil\)
and \(c(\delta_*+1)\ge K_0\); hence
\(L_0(\delta_*+1)=K_0(\delta_*+1)+c(\delta_*+1)(M-1)\ge K_0+K_0(M-1)=MK_0\).
The two-power refinement below shows that the tensor square of a
separator can be made exact with fewer copies than direct copying needs.

In the main construction direct copying would be too expensive. With
\(K_0=M_i\), it uses \(M_i^2\) source copies in each of the three
orientations of relation \(i\), so \(M_i^6\) in all, against \(M_i^3\)
output copies: a ratio \(M_i^3\) instead of \((2M_i-1)^3/M_i^3<8\). The copy
overhead would be \(3\sum_i\log M_i\), three times the contribution
\(\sum_i\log M_i\) of the separators to \(\log B_1\), instead of less than
\(3I\log2\) by~\eqref{eq:finite-cost-bounds}.

\subsection{A two-power refinement}

For the tensor square of the separator polynomial
\(\widetilde T(t)\) of Proposition~\ref{prop:separator-general}, an
average over a prime number of roots of unity removes all terms of
positive degree. Let \(\mathfrak p\) be a prime with
\(\mathfrak p>\delta_*\), \(\mathfrak p\equiv3\pmod 4\) and
\(\mathfrak p\nmid c\); such primes exist, since there are infinitely many
primes congruent to \(3\) modulo \(4\). Let \(\zeta_{\mathfrak p}\) be a
primitive \(\mathfrak p\)-th root of unity. Then
\begin{equation}\label{eq:two-power}
 \frac1{\mathfrak p}\sum_{\iota=0}^{\mathfrak p-1}
 \widetilde T(\zeta_{\mathfrak p}^\iota)^{\otimes2}=\widetilde T_0^{\otimes2}.
\end{equation}
Indeed, write
\(\widetilde T(t)=\sum_{\delta=0}^{\delta_*}t^{c\delta^2}\widetilde T_{c\delta^2}\),
where some coefficients may vanish. Then \(\widetilde T(t)^{\otimes2}\) is
the sum of the terms
\(t^{c(\delta_1^2+\delta_2^2)}\,\widetilde T_{c\delta_1^2}\otimes
\widetilde T_{c\delta_2^2}\) with \(0\le\delta_1,\delta_2\le\delta_*\),
and the average in~\eqref{eq:two-power} keeps exactly those with
\(\mathfrak p\mid c(\delta_1^2+\delta_2^2)\). Since \(\mathfrak p\nmid c\),
such a term has \(\delta_1^2+\delta_2^2\equiv0\pmod{\mathfrak p}\). If
\(\mathfrak p\nmid\delta_1\), then \((\delta_2\delta_1^{-1})^2\equiv-1
\pmod{\mathfrak p}\), which is impossible because \(-1\) is not a square
modulo a prime congruent to \(3\) modulo \(4\). Hence
\(\mathfrak p\mid\delta_1\), and then \(\mathfrak p\mid\delta_2\). As
\(0\le\delta_1,\delta_2\le\delta_*<\mathfrak p\), both vanish, and only
\(\widetilde T_0\otimes\widetilde T_0\) remains. Each evaluation
\(\widetilde T(\zeta_{\mathfrak p}^\iota)^{\otimes2}\) is the image of
\((L_0T)^{\otimes2}=L_0^2\,T^{\otimes2}\) under the tensor square of the
substitutions, so, as in the proof of
Lemma~\ref{lem:sparse-exactification},
\(\mathfrak p\,L_0^2\,T^{\otimes2}\restr
(\bigoplus_gT_gB_X(K_0))^{\otimes2}\).

For example, let \(M=K_0=8\) and \(c=1\). Then \(L_0=15\) and
\(\delta_*=7\), and the smallest admissible prime is \(\mathfrak p=11\).
The restriction above uses \(11\cdot15^2=2475\) copies of \(T^{\otimes2}\).
Direct copying applied to \(T^{\otimes2}\), which has \(M^2\) sectors and
needs a dot product of dimension \(K_0^2\), uses \((MK_0)^2=8^4=4096\)
copies. Lemma~\ref{lem:exactification} with \(k=2\) and degree
\(c\delta_*^2=49\) would need \(2\cdot49+1=99\) evaluations, that is,
\(99\cdot15^2=22275\) copies.
